% Please do not change %
\documentclass[11pt]{article} % font size
\usepackage[margin=1in]{geometry} % margins
\usepackage{amsmath,amsthm,amssymb} % for the  common "math symbols"
\usepackage{algorithm}
\usepackage{algpseudocode}
\usepackage[shortlabels]{enumitem} % for enumeration
\usepackage{booktabs} % if you need tables
\usepackage{listings}
\usepackage{color}
\usepackage{quiver}
\DeclareMathOperator*{\argmax}{argmax}
\DeclareMathOperator*{\argmin}{argmin}

% \theoremstyle{definition}
\newtheorem{theorem}{Theorem}[section]
\newtheorem{corollary}[theorem]{Corollary}
\newtheorem{lemma}[theorem]{Lemma}
\newtheorem{proposition}[theorem]{Proposition}
\newtheorem{conjecture}[theorem]{Conjecture}
\theoremstyle{definition}
\newtheorem{definition}[theorem]{Definition}
\theoremstyle{definition}
\newtheorem{fact}[theorem]{Fact}
\theoremstyle{definition}
\newtheorem{example}[theorem]{Example}

\lstset{
  frame=none,
  xleftmargin=2pt,
  stepnumber=1,
  numbers=left,
  numbersep=5pt,
  numberstyle=\ttfamily\tiny\color[gray]{0.3},
  belowcaptionskip=\bigskipamount,
  captionpos=b,
  escapeinside={*'}{'*},
  language=haskell,
  tabsize=2,
  emphstyle={\bf},
  commentstyle=\it,
  stringstyle=\mdseries\rmfamily,
  showspaces=false,
  keywordstyle=\bfseries\rmfamily,
  columns=flexible,
  basicstyle=\small\sffamily,
  showstringspaces=false,
  morecomment=[l]\%,
}
% Feel free to include additional packages (if needed), but make sure it does not override the margin/font requirements.
%
\setlength{\parindent}{0pt} % no indent for new paragraphs.
\setlength{\parskip}{5pt} % distance between paragraphs, which will be used when you have a blank line in your LaTeX code.


%%%%%%%%%%%%%%%%%%%%%%%%%%%%%%%%%%%%%%%%%%%%%%%%%%%%%%%%%%%%%%%%%%%%%%%%%%%%%%
\title{Categories, Proofs, and Processes: \\Assignment 3}
%
\author{Wira Azmoon Ahmad}
\date{11 Feb 2024}
%
\begin{document}
%
\maketitle

%You may use the usual \LaTeX\ environments to structure your answers:
%\begin{align}
%        \label{basegnn}
%        \mathbf{h}_u^{(t+1)} = \sigma \bigg( \mathbf{W}_\text{self}^{(t)} \mathbf{h}_u^{(t)}  + \mathbf{W}_\text{neigh}^{(t)} \sum_{v \in N(u)} \mathbf{h}_v^{(t)} + \mathbf{b}^{(t)} \bigg): 
%\end{align}


%\begin{table}[h]
%\caption{A simple table.}
%\centering
%\begin{tabular}[t]{lrrrr}
%\toprule
%Graphs & Can & Be  & Fun & Too \\ 
%\midrule 
%$G_1$ & $1$  & $2$ & $3$ & $4$ \\ 
%$G_2$ & $-1$  & $-2$ & $-3$ & $-4$ \\ 
%$G_3$ & $0.1$  & $0.2$ & $0.3$ & $0.4$ \\ 
%\bottomrule
%\end{tabular}
%\label{tab}
%\end{table}

\section*{Question 1}
\addtocounter{section}{1}

For each function $f:C\rightarrow D$ in \textbf{Set}, $S$ maps it to the functor
\[S\ f:\mathcal{C} \rightarrow \mathcal{D}\]
where $\mathcal{C},\mathcal{D}$ are discrete categories in \textbf{Cat} with $|C|,|D|$ objects respectively.

Then $\forall a\in C, b\in D, S\ a\in \text{Ob}(\mathcal{C}), S\ b\in \text{Ob}(\mathcal{D})$, the mapping is
\[S\ f::S\ a \mapsto S\ b \iff f::a\mapsto b\]

\begin{proposition}
    The functor $S:\textbf{Set}\rightarrow \textbf{Cat}$ is full and faithful.
\end{proposition}

\begin{proof}
To show faithfulness, let $f,g:C\rightarrow D$ be arbitrary functions in \textbf{Set} 
so $S\ f, S\ g: \mathcal{C} \rightarrow \mathcal{D}$.
Suppose $S\ f=S\ g$. Then since
\[S\ f::S\ a \mapsto S\ b \iff f::a\mapsto b\]
\[S\ g::S\ a \mapsto S\ b \iff g::a\mapsto b\]
we have 
\[f::a\mapsto b \iff S\ f=S\ g::S\ a \mapsto S\ b \iff g::a\mapsto b\]
and $f=g$. $S$ is faithful.

For fullness, let $F:\mathcal{C}\rightarrow \mathcal{D}$ be an arbitrary functor in \textbf{Cat}.
Then we can define $f:C\rightarrow D$ in \textbf{Set} 
such that $|\text{Ob}(\mathcal{C})|=|C|,|\text{Ob}(\mathcal{D})|=|D|$,
so $S\ f=F$. So $S$ is full.

\end{proof}

\clearpage
\section*{Question 2}
\addtocounter{section}{1}

\begin{definition}
The bivariant Hom-functor is given by
\[\mathcal{C}(\_,\_)(A,B):=\mathcal{C}(A,B)\]
\[\mathcal{C}(\_,\_)(A,f:B\rightarrow C):= \mathcal{C}(A,f),\quad
\mathcal{C}(\_,\_)(h:C\rightarrow B,A):= \mathcal{C}(h,A)\]
for all objects and arrows $A,B,C\in\text{Ob}(\mathcal{C}),
f\in\text{Ar}(\mathcal{C}),h\in\text{Ar}(\mathcal{C}^{\text{op}})$.
\end{definition}

\begin{proposition}
The bivariant Hom-functor is a functor.
\end{proposition}

\begin{proof}
Let $A,B,C,D\in \text{Ob}(\mathcal{C})=\text{Ob}(\mathcal{C}^{\text{op}})$ be arbitrary objects.

For any object pair $A,B$, the functor has a mapping to the set $\mathcal{C}(A,B)$.

Let $f:B\rightarrow C, g:C\rightarrow D$ be arbitrary arrows in $\mathcal{C}$.
\begin{equation}
\begin{aligned}
\mathcal{C}(\_,\_)(A,g\circ f)&= \mathcal{C}(A,g\circ f)\\
&= h \mapsto g\circ f\circ h\\
&= (h' \mapsto g\circ h')\circ (h'' \mapsto f\circ h'')\\
&= \mathcal{C}(A,g)\circ\mathcal{C}(A,f)\\
&= \mathcal{C}(\_,\_)(A,g)\circ\mathcal{C}(\_,\_)(A,f)\\
\end{aligned}
\end{equation}
Let $f':C\rightarrow B, g':D\rightarrow C$ be arbitrary arrows in $\mathcal{C}^{\text{op}}$.

\begin{equation}
\begin{aligned}
\mathcal{C}(\_,\_)(f'\circ g',A)&= \mathcal{C}(f'\circ g', A)\\
&= h \mapsto h\circ f'\circ g'\\
&= (h' \mapsto h'\circ g')\circ (h'' \mapsto h''\circ f')\\
&= \mathcal{C}(A,g')\circ\mathcal{C}(A,f')\\
&= \mathcal{C}(\_,\_)(A,g')\circ\mathcal{C}(\_,\_)(A,f')\\
\end{aligned}
\end{equation}
So the composition laws are satisfied.

For the identity laws,
\begin{equation}
\begin{aligned}
\mathcal{C}(\_,\_)(A,1_B)&= \mathcal{C}(A,1_B)\\
&= h\mapsto 1_B \circ h\\
&= h \mapsto h\\
&= 1_{\mathcal{C}(A,B)}\\
&= 1_{\mathcal{C}(\_,\_)(A,B)}\\
\end{aligned}
\end{equation}
\begin{equation}
\begin{aligned}
\mathcal{C}(\_,\_)(1_B,A)&= \mathcal{C}(1_B, A)\\
&= h\mapsto h\circ 1_B\\
&= h \mapsto h\\
&= 1_{\mathcal{C}(B,A)}\\
&= 1_{\mathcal{C}(\_,\_)(B,A)}\\
\end{aligned}
\end{equation}
Thus, the bivariant Hom-functor satisfies \textit{functoriality}.

\end{proof}

\section*{Question 3}
\addtocounter{section}{1}
\subsection*{(a)}
\begin{fact}
    Faithful functors do not, in general, reflect isomorphisms.
\end{fact}

\begin{example}
% https://q.uiver.app/#q=WzAsNCxbMCwwLCJBIl0sWzIsMCwiQiJdLFs2LDAsIkZcXCBBIl0sWzgsMCwiRlxcIEIiXSxbMCwxLCJmIiwwLHsib2Zmc2V0IjotMX1dLFsyLDMsIkZcXCBmIiwwLHsib2Zmc2V0IjotMX1dLFszLDIsIkZcXCBnPShGXFwgZileey0xfSIsMCx7Im9mZnNldCI6LTF9XSxbMSwwLCJnIiwwLHsib2Zmc2V0IjotMX1dXQ==
\[\begin{tikzcd}
	A \arrow["h", loop left] && B &&&& {F\ A} && {F\ B}
	\arrow["f", shift left, from=1-1, to=1-3]
	\arrow["{F\ f}", shift left, from=1-7, to=1-9]
	\arrow["{F\ g=(F\ f)^{-1}}", shift left, from=1-9, to=1-7]
	\arrow["g", shift left, from=1-3, to=1-1]
\end{tikzcd}\]
The diagram above shows two-object categories $\mathcal{C},\mathcal{D}$ respectively,
with the identities assumed but not drawn.
Note that $g\circ f=h\neq 1_A$, and $F\ h = 1_{F\ A}$.

$F\ f,F\ g$ in $\mathcal{D}$, are clearly mapped from $f,g$ in $\mathcal{C}$.
The functor $F:\mathcal{C}\rightarrow\mathcal{D}$ is faithful.
It can be noted it is not full since $F\ h=F\ 1_A= 1_{F\ A}$ but $h\neq 1_A$.
Fullness with faithfulness is the sufficient condition for reflecting isomorphism, proven in the next question.

The arrows $F\ f,F\ g$ are isomorphic to each other, but $f,g$ are not iso.
The functor does not reflect isomorphism.
\end{example}

\subsection*{(b)}
\begin{proposition}
    Every full and faithful functor reflects isomorphisms.
\end{proposition}
\begin{proof}
    Let the functor $F:\mathcal{C}\rightarrow \mathcal{D}$ be full and faithful.
    By faithfulness, for any arrow $g$ in $\mathcal{D}$, there exists an arrow $f$ in $\mathcal{C}$
    such that $F\ f = g$.

    Suppose we have two iso arrows $g,g^{-1}$ in $\mathcal{C}$.
    % where $F\ A,F\ B$ exist so
    Using faithfulness we find $f:A\rightarrow B,f':B' \rightarrow A'$ such that $F\ f=g,F\ f'=g^{-1}$.
    Since $g,g^{-1}$ are iso, $F\ A=F\ A',F\ B=F\ B'$ so 
    \[g:F\ A\rightarrow F\ B,\quad g^{-1}: F\ B \rightarrow F\ A\]
    where
    \[g^{-1}\circ g = 1_{F\ A},\quad g\circ g^{-1} = 1_{F\ B}\]
    By functoriality, we have from composition that 
    \[F(f\circ f')=F\ f\circ F\ f' = g \circ g^{-1}=1_{F\ B}\]
    \[F(f'\circ f)=F\ f'\circ F\ f = g^{-1} \circ g=1_{F\ A}\]
    We also have from the identity law that
    \[F(1_A)=1_{F\ A},\quad F(1_B)=1_{F\ B}\]
    By fullness,
    \[F(1_A)=F(f'\circ f)=1_{F\ A} \implies f'\circ f = 1_A\]
    \[F(1_b)=F(f\circ f')=1_{F\ B} \implies f\circ f' = 1_B\]
    So $f'=f^{-1}$, and $f,f^{-1}$ are iso.

    Given an iso arrow $F\ f$ in $\mathcal{D}$, we then have an iso arrow $f$ in $\mathcal{C}$.
    The full and faithful functor reflects isomorphisms.

\end{proof}

\subsection*{(c)}
\begin{proposition}
    Every equivalence preserves monics.
\end{proposition}

\begin{proof}
    Let $F:\mathcal{C}\rightarrow \mathcal{D}$ be an equivalence, 
    and let $f:B\rightarrow C$ be a monic arrow in $\mathcal{C}$.
    Then $F\ f:F\ B\rightarrow F\ C$.
    Let $F\ g, F\ h:A'\rightarrow F\ B$ be arbitrary arrows in $\mathcal{D}$.

    By equivalence, there exists an object $A$ of $\mathcal{C}$ such that $F(A)\cong A'$.

    By faithfulness, let the arrows $g,h:A\rightarrow B$ be in $\mathcal{C}$.

    Suppose $F\ f\circ F\ g = F\ f\circ F\ h$. Then
\begin{equation}
\begin{aligned}
    F\ f\circ F\ g = F\ f\circ F\ h &\implies F(f\circ g)=F(f\circ h) &&\quad \text{(by functoriality)}\\
    &\implies f\circ g = f\circ h &&\quad\text{(by fullness)}\\
    &\implies g=h &&\quad\text{(since $f$ is monic)}\\
    &\implies F\ g=F\ h
\end{aligned}
\end{equation}
    Then
    \[\forall F\ g, F\ h:F\ f\circ F\ g = F\ f\circ F\ h\implies F\ g= F\ h\]
    so the equivalence preserves monics.
\end{proof}

\clearpage
\section*{Question 4}
\addtocounter{section}{1}
\begin{proposition}
    If a category $\mathcal{C}$ has a terminal object and pullbacks, it has binary products and equalisers.
\end{proposition}
\begin{proof}
Let $T$ be the terminal object, and let $\tau_A:A\rightarrow T$ 
denote the unique arrow from an object $A$ to $T$.
% https://q.uiver.app/#q=WzAsNSxbMSwxLCJEIl0sWzAsMCwiRCciXSxbMiwxLCJCIl0sWzEsMiwiQSJdLFsyLDIsIlQiXSxbMSwwLCJoJyIsMix7InN0eWxlIjp7ImJvZHkiOnsibmFtZSI6ImRhc2hlZCJ9fX1dLFswLDIsInEiXSxbMSwyLCJxJyIsMCx7ImN1cnZlIjotMn1dLFswLDMsInAiLDJdLFsxLDMsInAnIiwyLHsiY3VydmUiOjJ9XSxbMyw0LCJcXHRhdV9BIiwyXSxbMiw0LCJcXHRhdV9CIl1d
\[\begin{tikzcd}
	{D'} \\
	& D & B \\
	& A & T
	\arrow["{h'}"', dashed, from=1-1, to=2-2]
	\arrow["q", from=2-2, to=2-3]
	\arrow["{q'}", curve={height=-12pt}, from=1-1, to=2-3]
	\arrow["p"', from=2-2, to=3-2]
	\arrow["{p'}"', curve={height=12pt}, from=1-1, to=3-2]
	\arrow["{\tau_A}"', from=3-2, to=3-3]
	\arrow["{\tau_B}", from=2-3, to=3-3]
\end{tikzcd}\]
The morphisms $A\xrightarrow[]{\tau_A}T \xleftarrow[]{\tau_B} B$ 
then have the pullback $A\xleftarrow[]{p} D\xrightarrow[]{q} B$.

For any morphisms $A\xleftarrow[]{p'} D'\xrightarrow[]{q'} B$, we have by composition,
\[\tau_A\circ p'=\tau_{D'}=\tau_B \circ q' \implies h':D'\rightarrow D \text{ is unique}\]
Then $D$ with morphisms $A\xleftarrow[]{p} D\xrightarrow[]{q} B$  is the product of $A$ and $B$,
since for all pairs $A\xleftarrow[]{p'} D'\xrightarrow[]{q'} B$, there is the unique morphism
$h'=\langle p', q'\rangle : D'\rightarrow D$ such that $p'=p\circ h'$ and $q'=q\circ h'$.

Using the binary product, we can then show equalisers.
Suppose we have the parallel morphisms $A\overset{f}{\underset{g}{\rightrightarrows}}B$.
As shown before, the product $A\times B$ 
with pullback morphisms $A\xleftarrow[]{\pi_1} D\xrightarrow[]{\pi_2} B$ exist.
The following diagram with $f$ commutes, and is similar if $g$ is used instead of $f$.
% https://q.uiver.app/#q=WzAsNCxbMiwyLCJBIl0sWzQsMCwiQiJdLFsyLDAsIkFcXHRpbWVzIEIiXSxbMCwwLCJBIl0sWzAsMSwiZiIsMl0sWzAsMiwiXFxsYW5nbGUxX0EsZlxccmFuZ2xlIiwyLHsib2Zmc2V0IjoxLCJzdHlsZSI6eyJib2R5Ijp7Im5hbWUiOiJkYXNoZWQifX19XSxbMiwxLCJcXHBpXzIiXSxbMiwzLCJcXHBpXzEiLDIseyJvZmZzZXQiOjF9XSxbMCwzLCIxX0EiXV0=
\[\begin{tikzcd}
	A && {A\times B} && B \\
	\\
	&& A
	\arrow["f"', from=3-3, to=1-5]
	\arrow["{\langle1_A,f\rangle}"', shift right, dashed, from=3-3, to=1-3]
	\arrow["{\pi_2}", from=1-3, to=1-5]
	\arrow["{\pi_1}"', shift right, from=1-3, to=1-1]
	\arrow["{1_A}", from=3-3, to=1-1]
\end{tikzcd}\]

Let $A\xleftarrow[]{p} E\xrightarrow[]{q} A$ be the pullback of 
$A\xrightarrow[]{\langle 1_A,f\rangle}A\times B\xleftarrow[]{\langle 1_A,g\rangle} A$.
% https://q.uiver.app/#q=WzAsNCxbMCwwLCJFIl0sWzAsMSwiQSJdLFsxLDEsIkFcXHRpbWVzIEIiXSxbMSwwLCJBIl0sWzAsMSwicCIsMl0sWzEsMiwiXFxsYW5nbGUgMV9BLCBmXFxyYW5nbGUiLDJdLFszLDIsIlxcbGFuZ2xlIDFfQSwgZ1xccmFuZ2xlIl0sWzAsMywicSJdXQ==
\[\begin{tikzcd}
	E & A \\
	A & {A\times B}
	\arrow["p"', from=1-1, to=2-1]
	\arrow["{\langle 1_A, f\rangle}"', from=2-1, to=2-2]
	\arrow["{\langle 1_A, g\rangle}", from=1-2, to=2-2]
	\arrow["q", from=1-1, to=1-2]
\end{tikzcd}\]
Then we show that $e=p=q:E\rightarrow A$ is the equaliser of $(f,g)$.
First, we show $p=q$ using the product,
\begin{equation}
\begin{aligned}
    p&=\pi_1 \circ \langle 1_A,f\rangle \circ p\\
    &= \pi_1 \circ \langle 1_A, g\rangle \circ q\\
    &= q 
\end{aligned}
\end{equation}
Then we have
\begin{equation}
\begin{aligned}
    f\circ e &= \pi_2\circ \langle 1_A, f\rangle \circ e\\
    &= \pi_2 \circ \langle 1_A, g\rangle \circ e\\
    &= g \circ e
\end{aligned}
\end{equation}
Suppose we then have a morphism $h:C\rightarrow A$ such that $f \circ h = g \circ h$. Then
\begin{equation}
\begin{aligned}
    \langle 1_A, f \rangle \circ h &= \langle 1_A \circ h, f \circ h\rangle\\
    &= \langle 1_A \circ h, g \circ h\rangle\\
    &= \langle 1_A, g\rangle \circ h
\end{aligned}
\end{equation}
Then by the pullback, there is a unique morphism $\hat{h}:C\rightarrow E$ 
such that $h=e\circ \hat{h}$. 
% https://q.uiver.app/#q=WzAsNCxbMSwwLCJBIl0sWzIsMCwiQiJdLFswLDAsIkUiXSxbMCwxLCJDIl0sWzAsMSwiZiIsMCx7Im9mZnNldCI6LTF9XSxbMCwxLCJnIiwyLHsib2Zmc2V0IjoxfV0sWzIsMCwiZSJdLFszLDAsImgiLDJdLFszLDIsIlxcaGF0e2h9IiwwLHsic3R5bGUiOnsiYm9keSI6eyJuYW1lIjoiZGFzaGVkIn19fV1d
\[\begin{tikzcd}
	E & A & B \\
	C
	\arrow["f", shift left, from=1-2, to=1-3]
	\arrow["g"', shift right, from=1-2, to=1-3]
	\arrow["e", from=1-1, to=1-2]
	\arrow["h"', from=2-1, to=1-2]
	\arrow["{\hat{h}}", dashed, from=2-1, to=1-1]
\end{tikzcd}\]

Thus, it is shown that the category $\mathcal{C}$ has binary products and equalisers.

\end{proof}

\section*{Question 5}
\addtocounter{section}{1}


\section*{Question 6}
\addtocounter{section}{1}



\end{document}
